\section{Checks and limiting cases}

\subsection{Boundary conditions}
At $x=0$: $G(0) = S(L)/S(L) = 1$, so $\mathcal{A}=1$, $\theta=0$ and
$u(0,t) = U_0\cos\omega t$, which satisfies \cref{eq:bc0}.
At $x=L$: $S(0) = 0$, so $\mathcal{A}(L)=0$ and $u(L,t)=0$, which satisfies \cref{eq:bcL}.

\subsection{Elastic limit, \texorpdfstring{$\eta \to 0$}{eta -> 0}}
With $\tau\to0$ we get $r\to1$, $\alpha\to0$ and $\beta\to\omega/c$. Then
$S(z)\to\sin(\omega z/c)$ is real, and
\begin{equation}
  u(x,t) \to U_0\,\frac{\sin\!\big(\omega(L-x)/c\big)}{\sin(\omega L/c)}\cos(\omega t).
\end{equation}
This is the classical undamped standing wave. It blows up at the resonances
$\omega_n = n\pi c/L$, where $\sin(\omega L/c) = 0$. With any damping
($\alpha>0$), \cref{eq:abs_sin} keeps the denominator bounded away from zero
and the resonant peaks become finite.

\subsection{Light damping, \texorpdfstring{$\omega\tau \ll 1$}{omega tau << 1}}
Expand $r = \sqrt{1+\omega^2\tau^2} \approx 1 + \tfrac12\omega^2\tau^2$. Then
$r - 1 \approx \tfrac12\omega^2\tau^2$, so to leading order
\begin{equation}
  \beta \approx \frac{\omega}{c}, \qquad
  \alpha \approx \frac{\omega}{c}\sqrt{\frac{\omega^2\tau^2}{4}} = \frac{\omega^2\tau}{2c}.
\end{equation}
The attenuation grows as $\omega^2$, a characteristic signature of
Kelvin--Voigt damping.

\subsection{Long or heavily damped rod, \texorpdfstring{$\alpha(L-x) \gg 1$}{alpha(L-x) >> 1}}
Write $\sin(kz) = \big(e^{ikz} - e^{-ikz}\big)/(2i)$. Since
$ikz = i\beta z + \alpha z$, the term $e^{ikz}$ grows like $e^{\alpha z}$ and
dominates when $\alpha z \gg 1$. Keeping only that term,
\begin{equation}
  G(x) \approx \frac{e^{ik(L-x)}}{e^{ikL}} = e^{-ikx} = e^{-\alpha x}\,e^{-i\beta x},
\end{equation}
so
\begin{equation}
  u(x,t) \approx U_0\, e^{-\alpha x}\cos(\omega t - \beta x).
\end{equation}
This is a single wave travelling to the right and decaying as it goes. The
wave reflected from $x = L$ has decayed to negligible size before it returns,
so the far boundary is effectively invisible.

\subsection{Travelling-wave interpretation}
In general, \cref{eq:X_solution} can be written as
\begin{equation}
  X(x) = \frac{U_0}{2i\sin(kL)}\Big[e^{ik(L-x)} - e^{-ik(L-x)}\Big].
\end{equation}
The first term is the incident wave travelling in $+x$. The second is its
reflection from the fixed end, inverted in sign so that the two cancel at
$x=L$.
